\section{Data Pipeline: Centroid Companion Files}
\label{sec:data-pipeline}

\subsection{File Format}

Geographic CVD requires a centroid companion file alongside the adjacency
file. The implemented loader derives the centroid filename from the adjacency
binary stem. For \texttt{\{stem\}.adj.bin}, it looks for
\texttt{\{stem\}\_centroids.json} in the same directory:

\begin{verbatim}
data/2020/
  north_carolina_adjacency_2020.adj.bin
  north_carolina_adjacency_2020_geoids.json
  north_carolina_adjacency_2020_centroids.json   <-- geographic CVD input
\end{verbatim}

The JSON structure is a single array of \texttt{[longitude, latitude]} pairs
in WGS84, one entry per census tract in the same index order as the
adjacency file:

\begin{verbatim}
{
  "centroids": [
    [-78.8740, 35.2271],   // tract 0
    [-80.8431, 35.2270],   // tract 1
    ...                    // one entry per tract
  ]
}
\end{verbatim}

\paragraph{Index alignment.}
The centroid file must have exactly $m$ entries, where $m$ is the number of
tracts in the adjacency file.
If $\mathrm{len}(\mathrm{centroids}) \ne m$, the loader returns an error:

\begin{verbatim}
ERROR: centroid file length mismatch: got 2659, expected 2660
       Check that centroids file was generated from the same
       adjacency file as the graph.
\end{verbatim}

\paragraph{Centroid definition.}
The Rust loader treats each entry as a WGS84 tract centroid aligned to
adjacency index order. The broader data package should record how those
centroids were derived: TIGER internal points, geometric polygon centroids, or
a population-weighted process. The algorithm consumes coordinate pairs; the
publication package must document their provenance.

\subsection{Fetching Centroids}
\label{subsec:fetch}

The CLI error path suggests the following user-facing fetch route:

\begin{verbatim}
bisect fetch --type centroids --states NC --year 2020
\end{verbatim}

This paper treats that command as the intended package workflow. A publication
package should verify the command, archive the generated
\texttt{\_centroids.json} files, and record hashes beside the adjacency and
GEOID files. To fetch centroids for all states once the fetch route is
available:

\begin{verbatim}
bisect fetch --type centroids --year 2020 --workers 8
\end{verbatim}

The fetch step should be idempotent: if the centroid file already exists and
its length matches the adjacency file, the fetch can be skipped.

\paragraph{CLI error when centroids absent.}
If \texttt{--cvd-metric geographic} is specified and the centroid companion
file is absent, the CLI returns:

\begin{verbatim}
ERROR: --cvd-metric geographic requires tract centroid data.
       Run: bisect fetch --type centroids
\end{verbatim}

No fallback to Phase~1 occurs silently: if the user requests geographic
metric and centroids are unavailable, the run fails with a clear diagnostic.

\subsection{LoadedGraph Integration}

The \texttt{LoadedGraph} struct gains one new field:

\begin{verbatim}
pub struct LoadedGraph {
    // ... existing fields ...
    /// WGS84 lon/lat tract centroid.
    /// Empty if the companion file is absent.
    pub tract_centroids: Vec<(f64, f64)>,
}
\end{verbatim}

Loading behaviour:
\begin{itemize}
\item After loading the adjacency file, derive the companion path from the
  adjacency binary stem and look for \texttt{\_centroids.json}.
\item If present and length-valid: load into \texttt{tract\_centroids}.
\item If absent: leave \texttt{tract\_centroids} empty (no error at load
  time; error is deferred to Phase~2 dispatch if \texttt{--cvd-metric geographic}
  is requested).
\item If present but length-mismatched: return error immediately at load time.
\end{itemize}

This design preserves full backward compatibility: all existing Phase~1
\cvd{} runs, all \metis{} runs, and all other structure-layer algorithms
are unaffected by the presence or absence of the centroid field.

\subsection{Consistency with Analysis Tools}

The Albers projection constants in \texttt{albers\_project}
(Definition~\ref{def:albers}) should be recorded beside the compactness
analysis configuration. The current transform uses central meridian
$-96^{\circ}$, parallels $29.5^{\circ}/45.5^{\circ}$, and
$R = 6{,}371{,}000$\,m.
A platform invariant test asserts that
$\mathrm{proj}(-96.0^{\circ}, 37.5^{\circ})$ returns $(|x| < 1{,}000\,\mathrm{m}, y > 0)$:
the central meridian projects to near-zero $x$ (within 1\,km) and positive $y$
(positive northing in the projection).
